\chapter{Construire un plan de connectivité}\label{ch:nw:build-a-connectivity-plan}

Une nouvelle société de trading va négocier sur trois plateformes~: des actions sur une bourse du New Jersey, des contrats à terme sur une bourse située hors de Chicago, et une plateforme crypto
dont le moteur d'appariement (matching engine) tourne dans une région cloud à Tokyo. Avant de signer quoi que ce soit, elle doit dire quelle latence elle obtiendra sur
chacune, ce qu'elle doit acheter, auprès de qui, à quelles conditions et ce que cela lui coûtera par an. Le dernier chapitre du livre répond par une
configuration unique, avec les composants du livre~: 1.02~million de dollars par an tout compris, trois des quatre cibles de latence atteintes, et une cible manquée
que corrige un \omterm{def:nw:private-connectivity-to-crypto-venues:endpoint}{point de terminaison privé} pour 123~dollars par an.

Chaque chapitre précédent a fourni une pièce. Le chapitre~10 a donné les sites et les \omterm{def:nw:the-north-american-map:floor}{planchers de latence}~; le chapitre~13 les routes en fibre et le chapitre~14
la radio~; les chapitres~2 à~8 la cage, les cartes, le chemin logiciel et les serveurs~; le chapitre~9 les prix de la \omterm{def:nw:colocation-products-and-how-they-are-sold:colo}{colocation}~; le chapitre~15 les circuits et les niveaux de
service~; les chapitres~16 à~18 le cloud~; le chapitre~27 la carte des fournisseurs~; le chapitre~28 la reprise. Ce chapitre les réunit dans
\texttt{firm.connplan}, le couronnement du livre, et l'exécute sur un plan dont les prix sont des lignes datées issues des grilles tarifaires des plateformes elles-mêmes.

\section{Exigences~: stratégies, plateformes et cibles de latence}

\begin{definition}[Plan de connectivité, nomenclature]\label{def:nw:build-a-connectivity-plan:plan}
Un \emph{plan de connectivité}\index{plan de connectivité} est un document de la société, conservé sous forme de données, qui indique, pour chaque plateforme sur laquelle elle négocie,
la latence dont elle a besoin et la manière dont elle l'obtiendra~: les sites, les routes, les liaisons et les services qu'elle achètera, leurs niveaux de service et
leurs dispositifs de reprise, ainsi que leur coût. Sa \emph{nomenclature}\index{nomenclature} (bill of materials) est la liste détaillée de tout ce qu'il faut acheter
avec la quantité, le fournisseur, le prix et la source datée de chaque prix.
\end{definition}

Les stratégies de la société fixent les cibles. Sa stratégie sur actions réagit au carnet d'ordres propre à la plateforme d'actions~; elle se place donc en \omterm{def:nw:colocation-products-and-how-they-are-sold:colo}{colocation} à Mahwah et
a besoin d'un temps de bout en bout inférieur à \qty{40}{\micro\second} au 99\up{e}~centile. Sa stratégie sur contrats à terme se couvre sur le CME~: elle a besoin des cours d'Aurora
à Mahwah en moins de \qty{4.2}{\milli\second} et de ses ordres à Aurora en moins de \qty{8}{\milli\second}. Sa stratégie crypto négocie sur une
plateforme située dans la région Tokyo d'AWS, depuis des instances placées dans cette région, et a besoin d'allers-retours vers la plateforme inférieurs à une milliseconde au 99\up{e}~centile
(toutes les cibles sont celles de la société et sont ici des hypothèses).

\omcode{code/networks/29-build-a-connectivity-plan/python/nw_plan.py}{17}{36}{Le plan sous forme de données~: plateformes avec leur accès et leurs cibles, et périmètre des liaisons tarifées.}\label{lst:nw:build-a-connectivity-plan:plan}

\section{La table des latences}

Chaque type d'accès tire sa latence d'un composant~: un chemin colocalisé, issu de la cage du chapitre~2 et du chemin logiciel du chapitre~5
(\texttt{firm.cagenet}, \texttt{firm.wirepath}, les étapes mesurées du livre~13)~; une route radio, à partir de son chiffre publié~; une route en fibre
à partir de la géodésique et d'un \omterm{def:nw:the-north-american-map:factor}{facteur de tracé} (\texttt{firm.fibreroute})~; un chemin cloud, issu des chemins privés du chapitre~18
(\texttt{firm.privlink}).

\omcode{code/firm/connplan/firm_connplan.py}{100}{111}{Une latence par plateforme, issue du composant qui modélise son accès.}\label{lst:nw:build-a-connectivity-plan:latency}

\begin{table}[ht]
\centering\small
\begin{tabular}{@{}llrrl@{}}
\toprule
Plateforme & Mesure & Obtenue & Cible & \\
\midrule
Actions NYSE (Mahwah) & bout en bout, 99\up{e}~centile & \qty{35.6}{\micro\second} & \qty{40}{\micro\second} & atteinte \\
Données CME, Aurora vers Mahwah & un sens (radio) & \qty{3.986}{\milli\second} & \qty{4.2}{\milli\second} & atteinte \\
Ordres CME, Mahwah vers Aurora & un sens (fibre à 1.3) & \qty{7.48}{\milli\second} & \qty{8}{\milli\second} & atteinte \\
Binance (AWS Tokyo) & aller-retour, 99\up{e}~centile & \qty{2.40}{\milli\second} & \qty{1}{\milli\second} & manquée \\
\bottomrule
\end{tabular}
\caption{La table des latences du plan. Le chiffre radio est la route Aurora--Mahwah publiée la plus rapide (chapitre~10)~; la grille tarifaire de NYSE
ne donne aucune latence pour ses données CME sans fil. Le facteur de la route en fibre et le temps du chemin cloud jusqu'à la périphérie sont des hypothèses. Données~:
\texttt{nw\_plan.report()}.}\label{tab:nw:build-a-connectivity-plan:latency}
\end{table}

\begin{omfigure}
\begin{tikzpicture}
\begin{axis}[ybar,width=10cm,height=5.4cm,bar width=16pt,ymin=0,ymax=280,xtick=data,
  xticklabels={NYSE,CME données,CME ordres,Binance},
  ylabel={latence obtenue, \% de la cible},nodes near coords,
  every node near coord/.append style={font=\scriptsize,anchor=north,yshift=-2pt},
  nodes near coords={\pgfmathprintnumber[fixed,fixed zerofill,precision=1]{\pgfplotspointmeta}},
  tick label style={font=\footnotesize},label style={font=\small},grid=major,grid style={gray!25},table/col sep=comma]
\draw[omThm,thick,dashed] (axis cs:-0.5,100) -- (axis cs:3.5,100);
\addplot[fill=omDef!55,draw=omDef] table[x=x,y=share]{figdata/networks/29-build-a-connectivity-plan/latency.csv};
\end{axis}
\end{tikzpicture}

\omcaption{La latence obtenue de chaque plateforme, en part de sa cible~; la ligne en pointillé représente la cible. Trois plateformes s'y trouvent~; la
plateforme crypto de Tokyo, atteinte par la périphérie publique d'internet, est à 240~\%. Données~: \texttt{fig\_plan.py},
\texttt{nw\_plan.report()}.}\label{fig:nw:build-a-connectivity-plan:latency}
\end{omfigure}

La cible manquée est celle que le plan peut corriger à peu de frais. La configuration donne à la plateforme crypto les trois chemins du chapitre~18 (une
hypothèse sur ce qu'elle propose)~: son point de terminaison public à travers un \omterm{def:nw:where-the-crypto-venues-live:as}{réseau de périphérie} (edge), et un \omterm{def:nw:private-connectivity-to-crypto-venues:endpoint}{point de terminaison privé} dans la zone de la société ou dans une autre zone. Le plan les essaie tous et retient ceux qui atteignent la cible, classés par
coût.

\omcode{code/firm/connplan/firm_connplan.py}{154}{169}{La modification la moins chère qui atteint une cible cloud manquée~: chaque chemin proposé, essayé et chiffré.}\label{lst:nw:build-a-connectivity-plan:fix}

Un \omterm{def:nw:private-connectivity-to-crypto-venues:endpoint}{point de terminaison privé} dans la zone de l'instance ramène l'aller-retour au 99\up{e}~centile de \qty{2.40}{\milli\second} à
\qty{0.44}{\milli\second}, pour 122.64~dollars par an de durée d'utilisation du point de terminaison (au prix du chapitre~18~; sa facturation du trafic s'y ajoute avec le volume)~;
un point dans une autre zone atteint aussi la cible, à \qty{0.80}{\milli\second}, pour le même prix fixe, plus le transfert inter-zones. La correction consiste à
aligner le point de terminaison sur l'identifiant de zone de l'instance, comme le montre la leçon du chapitre~16.

\section{La nomenclature}

\begin{dated}{2026-09}{Prix du plan, issus des grilles tarifaires publiées des plateformes et du cloud}\label{dat:nw:build-a-connectivity-plan:prices}
NYSE (grille du 14 septembre 2026)~: baie dédiée 5\,000 dollars à l'ouverture et 1\,200 dollars par kW et par mois pour 4 à 8 kW~; connexion réseau 10 Gb LCN et NMS
15\,000 dollars à l'ouverture et 24\,000 par mois (22\,000 dans la grille d'octobre 2025 de NYSE, telle que la cite le dépôt de Nasdaq)~; \omterm{def:nw:colocation-products-and-how-they-are-sold:mmr}{interconnexion}
(cross-connect) 500 et 600 par mois~; données CME sans fil 5\,000 et 6\,000 par mois~; flux de données CME sur le réseau IP 3\,000 par mois. La
connexion Ultra 10 Gb comparable de Nasdaq~: 18\,500 dollars par mois à partir de janvier 2026. AWS Tokyo (grille tarifaire du 25 septembre 2026)~: c7i.4xlarge
0.8988~dollar de l'heure.
\end{dated}

\begin{table}[ht]
\centering\footnotesize
\begin{tabular}{@{}lrrrl@{}}
\toprule
Élément & Qté & Mensuel (dollars) & Annuel (dollars) & Source \\
\midrule
Baie dédiée NYSE (frais d'ouverture) & 1 & -- & 1\,667 & 29~:F1 \\
Alimentation de la baie NYSE, 4 à 8 kW & 8 kW & 9\,600 & 115\,200 & 29~:F1 \\
\omterm{def:nw:colocation-products-and-how-they-are-sold:mmr}{Interconnexions} en fibre NYSE & 4 & 2\,400 & 29\,467 & 29~:F1 \\
Connexion NYSE LCN et NMS 10 Gb & 1 & 24\,000 & 293\,000 & 29~:F1 \\
Réseau IP NYSE 1 Gb (secours) & 1 & 2\,500 & 30\,833 & 15~:F4 \\
Longueur d'onde Mahwah--Aurora (ordres CME) & 1 & 20\,000 & 243\,333 & hypothèse \\
Données CME sans fil NYSE & 1 & 6\,000 & 73\,667 & 29~:F1 \\
Flux de données CME sur le réseau IP & 1 & 3\,000 & 36\,000 & 29~:F1 \\
AWS Tokyo c7i.4xlarge (730 h) & 2 & 1\,312 & 15\,747 & 29~:F2 \\
\bottomrule
\end{tabular}
\caption{La \omterm{def:nw:build-a-connectivity-plan:plan}{nomenclature} du plan~; les frais ponctuels sont amortis sur 36 mois dans la colonne annuelle. La longueur d'onde
n'a pas de prix publié~: le plan retient un prix supposé, que son propre contrôle signale. Données~: \texttt{code/firm/connplan/data/cost\_table.csv}.}\label{tab:nw:build-a-connectivity-plan:bom}
\end{table}

Deux contrôles sont exécutés avant que le plan ne soit jugé fiable~: chaque ligne de prix doit nommer sa source de grand livre, et aucune ne peut dater de plus de 60 jours. Le jour
de la rédaction, une ligne échoue au premier contrôle, la longueur d'onde Mahwah--Aurora, dont les 20\,000~USD par mois sont une hypothèse que la société doit remplacer
par la cotation d'un opérateur~; aucune n'échoue au second. Le tableau est aussi l'export du plan~: \texttt{export\_cost\_table} l'écrit avec ses
sources et dates pour les responsables budgétaires de la société (le livre~16 le lit).

\section{Contrats et niveaux de service}

L'accès du plan à la plateforme d'actions repose sur une connexion de 10~Gb, avec un circuit IP de 1~Gb en secours. Chaque circuit tombant en panne
une fois sur 4\,000 heures et réparé en quatre heures (hypothèses), une seule connexion est disponible 99.900~\% du temps~; avec le secours et un
événement de mode commun (une gaine ou une salle partagée) deux fois en dix ans pendant huit heures, 99.982~\% (\texttt{firm.slamodel}, chapitre~15). Une seconde
connexion de 10~Gb à côté du secours ajouterait 366\,250~USD par an au budget, avec la quote-part de reprise~: le prix du maintien de la pleine bande passante pendant une panne.
Le barème de frais lu pour le plan ne mentionne aucun
crédit de niveau de service pour ces circuits~; ceux-ci relèvent du contrat.

\section{Le budget annuel et ses sensibilités}

\begin{omfigure}
\begin{tikzpicture}
\begin{axis}[xbar,width=10cm,height=5cm,bar width=9pt,xmin=0,xmax=650,ytick=data,y dir=reverse,
  yticklabels={colocation,connectivité,données de marché,cloud,reprise après sinistre},
  xlabel={milliers d'USD par an},nodes near coords,every node near coord/.append style={font=\scriptsize},
  nodes near coords={\pgfmathprintnumber[fixed,fixed zerofill,precision=1]{\pgfplotspointmeta}},
  tick label style={font=\footnotesize},label style={font=\small},grid=major,grid style={gray!25},table/col sep=comma]
\addplot[fill=omThm!45,draw=omThm] table[y=x,x=kusd]{figdata/networks/29-build-a-connectivity-plan/budget.csv};
\end{axis}
\end{tikzpicture}

\omcaption{Le budget annuel du plan par catégorie, 1.02~million d'USD au total. La connectivité domine~; la reprise après sinistre vaut un quart de la
\omterm{def:nw:colocation-products-and-how-they-are-sold:colo}{colocation} sur site et de la connectivité (conception tiède du chapitre~28, une hypothèse). Données~: \texttt{fig\_plan.py},
\texttt{nw\_plan.report()}.}\label{fig:nw:build-a-connectivity-plan:budget}
\end{omfigure}

Le plan coûte 68\,812~USD par mois, plus 39\,500 de frais ponctuels~: 1\,017\,289~USD par an, avec les frais amortis sur trois ans et le
site de reprise à un quart des dépenses sur site. La connectivité en représente 56~\%, et une seule connexion, la liaison LCN et NMS de 10~Gb, en représente 29~\%
à elle seule. Les sensibilités montrent où le budget bouge~: une seconde connexion LCN ajoute 366\,250~USD par an~; quatre kW de puissance en plus, 72\,000~USD~; une
longueur d'onde facturée 50~\% au-dessus de l'hypothèse, soit 152\,083~USD~; doubler les instances de Tokyo n'ajoute que 15\,747~USD~; retirer le circuit de secours économise 38\,542~USD
et fait perdre la disponibilité indiquée plus haut. La branche cloud, celle qui a manqué sa cible, est la moins chère à corriger et à faire grandir.

\begin{method}[Construire et tenir à jour un plan de connectivité]\label{met:nw:build-a-connectivity-plan:method}
\begin{enumerate}
\item Écrire les exigences sous forme de données~: chaque plateforme, ce que chaque stratégie en attend, et la mesure (centile, aller simple ou aller-retour).
\item Calculer chaque latence avec le composant qui modélise son accès~; comparer à la cible et signaler chaque écart.
\item Pour chaque écart, essayer les alternatives proposées et classer par coût celles qui atteignent la cible.
\item Chiffrer chaque ligne à partir d'une donnée datée et sourcée~; signaler les lignes sans source et périmées, et remplacer les hypothèses par des cotations avant la signature.
\item Ajouter les niveaux de service et la conception de reprise~; calculer le budget et ses sensibilités~; exporter le tableau des coûts.
\item Relancer le plan à chaque changement de prix, de plateforme ou de stratégie~: le plan est un programme, non un document.
\end{enumerate}
\end{method}

\section{Tutoriel~: trois plateformes, un seul budget}\label{tut:nw:build-a-connectivity-plan:tut}

\begin{tutorial}
\textbf{Objectif.} Exécuter tout le plan à partir d'une seule configuration et modifier une exigence.
\textbf{État final~:} les \cref{tab:nw:build-a-connectivity-plan:latency,tab:nw:build-a-connectivity-plan:bom} et
la \cref{fig:nw:build-a-connectivity-plan:budget}.
\begin{tutsteps}
\item \textbf{Configuration.} \texttt{nw\_plan.VENUES} et \texttt{LINES} (le \cref{lst:nw:build-a-connectivity-plan:plan}) avec
  \texttt{firm\_connplan.load\_prices}.
\item \textbf{Latence.} \texttt{latency\_table} (le \cref{lst:nw:build-a-connectivity-plan:latency})~; \texttt{cheapest\_fixes}
  (le \cref{lst:nw:build-a-connectivity-plan:fix}).
\item \textbf{Coûts.} \texttt{bom}, \texttt{budget}, \texttt{check\_prices}, \texttt{nw\_plan.sensitivities} et
  \texttt{export\_cost\_table}.
\end{tutsteps}
\textbf{Pour aller plus loin.} Resserrer la cible sur actions à \qty{20}{\micro\second} et laisser le plan proposer le chemin matériel du chapitre~7~;
déplacer la stratégie sur contrats à terme vers Aurora et comparer les budgets.
\end{tutorial}

\section{Construction~: le plan de connectivité}\label{bld:nw:build-a-connectivity-plan:connplan}

\begin{build}
\bfield{Rôle} Le \omterm{def:nw:build-a-connectivity-plan:plan}{plan de connectivité} de la société sous forme de programme~: tableau des latences, \omterm{def:nw:build-a-connectivity-plan:plan}{nomenclature}, contrats, ligne de reprise, budget,
sensibilités, corrections, contrôles et tableau des coûts exporté. Le point d'orgue du livre.
\bfield{Interface} \texttt{firm\_connplan}~: \texttt{PriceRow}, \texttt{load\_prices}, \texttt{Venue}, \texttt{Line}, \texttt{Plan},
\texttt{latency\_us}, \texttt{latency\_table}, \texttt{bom}, \texttt{budget}, \texttt{check\_prices}, \texttt{cheapest\_fixes},
\texttt{availability}, \texttt{export\_cost\_table}~; les colonnes du tableau exporté dans son \texttt{README.md}.
\bfield{Règles} Chaque ligne de prix nomme sa source de grand livre et sa date, ou est signalée~; les latences viennent des composants du livre~; les cibles et
hypothèses sont énoncées dans la configuration.
\bfield{Tests d'acceptation} \texttt{code/firm/connplan/tests/}~: lignes de prix, un budget calculé à la main, les contrôles et la correction sur un petit plan, un tracé de fibre
comparé à son \omterm{def:nw:the-north-american-map:floor}{plancher de latence}, et le tableau exporté.
\bfield{Pour aller plus loin} Chaque type d'accès pour chaque plateforme, avec ses solutions de rechange chiffrées~; conversion de devises avec des taux datés~; le chemin matériel
comme option pour les plateformes en \omterm{def:nw:colocation-products-and-how-they-are-sold:colo}{colocation}.
\end{build}

\begin{omsources}
\item NYSE, Connectivity Fees and Charges (septembre 2026)~; Nasdaq, SR-NASDAQ-2025-089~; liste de prix AWS EC2 (Tokyo, septembre 2026).
\item Les chapitres~2 à~28 de ce livre et leurs grands livres.
\end{omsources}

\section{Exercices}

\begin{exercise}[$\star$]\label{exo:nw:build-a-connectivity-plan:1}
Définir ce qu'est une \omterm{def:nw:build-a-connectivity-plan:plan}{nomenclature} et préciser ce que chacune de ses lignes doit porter dans le plan de ce livre.
\end{exercise}

\begin{exercise}[$\star$]\label{exo:nw:build-a-connectivity-plan:2}
Calculer le coût mensuel d'une baie dédiée de 8~kW à Mahwah selon le barème de septembre 2026.
\end{exercise}

\begin{exercise}[$\star$]\label{exo:nw:build-a-connectivity-plan:3}
Expliquer pourquoi le plan signale la longueur d'onde Mahwah--Aurora.
\end{exercise}

\begin{exercise}[$\star\star$]\label{exo:nw:build-a-connectivity-plan:4}
Calculer la part du budget annuel que représente la connexion LCN et NMS de 10~Gb, et la hausse de son prix mensuel entre les barèmes
de NYSE d'octobre 2025 et de septembre 2026.
\end{exercise}

\begin{exercise}[$\star\star$]\label{exo:nw:build-a-connectivity-plan:5}
Expliquer pourquoi un \omterm{def:nw:private-connectivity-to-crypto-venues:endpoint}{point de terminaison privé} situé dans une autre zone atteint la cible de Tokyo, mais coûte encore plus cher que celui de la zone de l'instance.
\end{exercise}

\begin{exercise}[$\star\star$]\label{exo:nw:build-a-connectivity-plan:6}
Quelle modification unique fait le plus varier le budget, et pourquoi n'est-elle pas évidemment rentable~?
\end{exercise}

\begin{exercise}[$\star\star\star$]\label{exo:nw:build-a-connectivity-plan:7}
\emph{Programmation.} Avec \texttt{firm\_connplan}, quelle latence unidirectionnelle l'itinéraire d'ordres vers le CME atteint-il avec un \omterm{def:nw:the-north-american-map:factor}{facteur de tracé} de 1.2 au lieu de 1.3~?
\end{exercise}

\begin{exercise}[$\star\star\star$]\label{exo:nw:build-a-connectivity-plan:8}
\emph{Trouver la faille.} «~Notre plan coûte 1.02~million de dollars par an~; nous pouvons signer aujourd'hui.~»
\end{exercise}

\section{Problème~: trois plateformes, un seul budget}

\begin{problem}[{Problème du week-end --- le plan de synthèse}]\label{pb:nw:build-a-connectivity-plan:1}
La société de trading du chapitre doit présenter son plan à son conseil d'administration. Utilisez \texttt{firm.connplan} et la configuration du chapitre.

\medskip\noindent\textbf{Partie I --- Latence.}
\begin{enumerate}
\item Qu'obtient chaque plateforme par rapport à sa cible~?
\item Quels composants ont calculé chaque chiffre~?
\item Lesquels des chiffres sont publiés et lesquels sont modélisés~?
\item Qu'est-ce qui n'est pas atteint, de combien et pourquoi~?
\end{enumerate}

\medskip\noindent\textbf{Partie II --- La correction.}
\begin{enumerate}[resume]
\item Quelles options atteignent la cible manquée, et à quelle latence~?
\item Que coûte la solution la moins chère par an~?
\item Que doit vérifier la société dans son compte cloud pour que la correction fonctionne~?
\item Pourquoi la correction est-elle si peu coûteuse par rapport au reste du plan~?
\end{enumerate}

\medskip\noindent\textbf{Partie III --- Coûts.}
\begin{enumerate}[resume]
\item Quels sont les coûts mensuels, ponctuels et annuels~?
\item Comment se répartit le coût annuel par catégorie~?
\item Quelle ligne échoue aux contrôles, et que doit faire la société~?
\item Que montrent les sensibilités~?
\end{enumerate}

\medskip\noindent\textbf{Partie IV --- Le verdict.}
\begin{enumerate}[resume]
\item Énoncer le \emph{résultat nommé}~: le coût annuel total tout compris de l'entrée dans les trois plateformes, la latence atteinte
  sur chacune par rapport à sa cible, et le changement le moins coûteux qui permet d'atteindre l'unique cible manquée.
\item Qu'apporte le circuit de secours~?
\item Que changerait le transfert de la stratégie sur contrats à terme vers Aurora~?
\item Que doit demander le conseil avant d'approuver~?
\item À quelle fréquence faut-il relancer le plan~?
\item Que retient le livre~16 du plan~?
\item Quelle partie du plan faudrait-il vérifier en premier si une cible de latence était manquée en production~?
\item En une phrase~: à quoi sert un \omterm{def:nw:build-a-connectivity-plan:plan}{plan de connectivité}~?
\end{enumerate}
\end{problem}

\section{Questions d'entretien}

\begin{interviewq}[$\star$ \iqroles{développeur}]\label{iq:nw:build-a-connectivity-plan:1}
Que contient le \omterm{def:nw:build-a-connectivity-plan:plan}{plan de connectivité} d'une société de trading~?
\end{interviewq}

\begin{interviewq}[$\star\star$ \iqroles{développeur}]\label{iq:nw:build-a-connectivity-plan:2}
Comment estimeriez-vous la latence entre le New Jersey et une bourse de contrats à terme près de Chicago avant d'acheter quoi que ce soit~?
\end{interviewq}

\begin{interviewq}[$\star\star$ \iqroles{développeur, chercheur}]\label{iq:nw:build-a-connectivity-plan:3}
Votre plan manque une cible de latence. Comment décidez-vous de ce qu'il faut changer~?
\end{interviewq}

\begin{interviewq}[$\star\star$ \iqroles{développeur}]\label{iq:nw:build-a-connectivity-plan:4}
Comment gardez-vous les prix d'un plan fiables dans le temps~?
\end{interviewq}

\begin{interviewq}[$\star\star$ \iqroles{chercheur}]\label{iq:nw:build-a-connectivity-plan:5}
Le conseil demande si la seconde connexion de 10~Gb vaut 366\,000 dollars par an. Que répondez-vous~?
\end{interviewq}

\begin{interviewq}[$\star\star\star$ \iqroles{développeur, chercheur}]\label{iq:nw:build-a-connectivity-plan:6}
Concevez la connectivité d'une nouvelle société de trading d'actions américaines, de contrats à terme du CME et d'une plateforme crypto à Tokyo, depuis les besoins jusqu'au budget.
\end{interviewq}
